% LaTeX Curriculum Vitae
% Author: Xovee Xu (https://xovee.cn)
% Last Update: July 28, 2020
% Latest Version (consider give me a star!): https://github.com/Xovee/latex-cv
% Overleaf Link: https://www.overleaf.com/latex/templates/xovees-cv-template/rrsmqwhbygcf
% Current Version: v0.95
% LICENSE: MIT


\documentclass{article}

\usepackage[utf8]{inputenc}
\usepackage[full]{textcomp}
\usepackage{CJKutf8}
\usepackage[lf]{ebgaramond}
% \usepackage[scaled,swashQ]{garamondx}
\usepackage[T1]{fontenc}

\usepackage{enumitem}
\usepackage[a4paper,left=.9in, right=.9in, top=1.5in, bottom=1.5in]{geometry}
\usepackage{url}
\usepackage[dvipsnames]{xcolor}
\usepackage{amsmath}

% package settings
\usepackage[
    hidelinks,
    pdfnewwindow=true,
    pdfauthor={Shih-Yang Lin},
    pdftitle={Curriculum Vitae of Shih-Yang Lin},
]{hyperref}

\pagestyle{headings}
\markright{\textbf{Shih-Yang Lin}}

\setlength\parindent{2em}

\thispagestyle{empty}

% define cv section
\newcommand{\cvsection}[1]{\section*{\rmfamily#1}}
\newcommand{\cvsubsection}[1]{\subsection*{\rmfamily\hspace{1.6em}#1}}

% begin
\begin{document}

% name
\begin{center}
    \Huge{
    \rmfamily
    \textbf{Shih-Yang Lin}}
\end{center}
\vspace{20pt}


\setlength{\parskip}{1pt}
\renewcommand{\arraystretch}{1.25}

% Contact Information

\noindent Department of Economics \hfill 0982-031-228

\noindent National Taiwan University (NTU)\hfill \href{mailto:d07323003@ntu.edu.tw}{\texttt{d07323003@ntu.edu.tw}}

\noindent No.1, Sec 4, Roosevelt Rd., Taipei 10617, Taiwan\hfill %\url{https://xovee.cn}


\setlength{\parskip}{3pt}




% Summary / Statement
%\cvsection{Summary / Statement}
%\indent

%\hangindent=2emI have rich experience and strong skills in ... I want ...

\cvsection{Research of Interests}
\indent 

\textbf{Macroeconomics, Econometrics}


% Education
\cvsection{Education}
\indent 

\textbf{National Taiwan University (NTU)}, Taipei, Taiwan

\hspace{2em}Ph.D. Student, Economics, 2018 to 2024 (expc.)

\hspace{2em}M.A., Economics, 2016

\textbf{National Sun Yat-Sen University (NSYSU)}, Kaohsiung, Taiwan

\hspace{2em}B.A., Political Economy, 2008


% below is another kind of style
%\cvsection{Education (Another Style)}
%\indent 

%Ph.D., Computer Science, University of Electronic Science and Technology of China, 2021 to 2025 (expc.)

%M.S., Software Engineering, University of Electronic Science and Technology of China, 2021

%B.S., Software Engineering, University of Electronic Science and Technology of China, 2018


% academic employment
\cvsection{Academic Employment}
\indent

\textbf{Academia Sinica}, Taipei, Taiwan

\hspace{2em}Research Assistant, 2016 - 2018


%\cvsection{Academic Employment (Another Style)}
%\indent

%Dean, Dept. of Computer Science, Standard University, CA, USA, 2080 -

%Professor, Standard University, CA, USA, 2077 -

%Associated Professor, Standard University, CA, USA, 2066 - 2077

%Assistant Professor, Standard University, CA, USA, 2060 - 2066

%Postdoc, University of Electronic Science and Technology of China, Chengdu, China, 2059 - 2060

% distinction, award, honor, and fellowship
\cvsection{Grants and Awards}
\indent

Doctoral Dissertation Fellowships of the Humanities and Social Sciences, National Science and Technology Council, 2023.

Asia Bank Scholarship, National Taiwan University, 2022.

Liu Ta-Chung Scholarship, National Taiwan University, 2020.

Harold H.C. Han Scholarship, National Taiwan University, 2019.

Chu, Cyrus C. Y. and Chen, Tain-Jy Scholarship, National Taiwan University, 2018.

% publication
\section*{Publication}

\cvsubsection{Refereed Journal Articles}

\begin{enumerate}[resume]
    \item \textbf{“Estimating the Willingness to Pay for Voting when Absentee Voting is not Allowed.”} Coauthored with C.Y. Cyrus Chu and Wen-Jen Tsay. 2021. \textit{Social Science Quarterly}, 102: (4), pp. 1380-1393,\\ doi:\href{https://doi.org/10.1111/ssqu.13014}{10.1111/ssqu.13014}.
\end{enumerate}


\cvsubsection{Non-refereed Journal Articles}

\begin{enumerate}[resume]
    \item \textbf{“Constructing the DSGE-VAR Model-Taiwan's Medium-to Long-Term Economic Growth Rate Forecast.”} Coauthored with Yi-Chan Tsai, Ruey Yau, Kuo-Hsan Chin, 2022. \textit{Quarterly Bulletin, Central Bank of the Republic of China (Taiwan)}, 44: (2), pp. 27-70.
\end{enumerate}

\section*{Working Paper}

\begin{enumerate}[resume]
    \item \textbf{“Financial Frictions and Uncertainty Shocks.”} Coauthored with Yi-Chan Tsai and Po-Yuan Wang. 
    \begin{itemize}
        \item[] This paper investigates the role of credit constraints in the propagation of uncertainty shocks within a business cycle framework. Traditional business cycle models struggle to account for observed declines in output, consumption, investment, and labor hours in response to heightened uncertainty. To address this gap, we introduce a collateral-based credit constraint for both impatient households and entrepreneurs, with borrowing limitations connected to the value of their collateral assets. As uncertainty escalates, households and entrepreneurs demand a higher risk premium for collateral, leading to a decrease in its overall demand. This reduced demand, in turn, prompts households to curtail their labor supply, leading to a decrease in output. Our study emphasizes that collateral adjustments following uncertainty shocks can induce macroeconomic co-movements in a real business cycle model, even in the absence of nominal rigidities. Furthermore, we discover that a lower loan-to-value (LTV) ratio can help alleviate the adverse impacts of uncertainty shocks. This research offers a new perspective on how financial constraints shape macroeconomic dynamics in times of heightened uncertainty, providing valuable insights for policymakers and economists.
    \end{itemize}
    \item \textbf{“The Labor Supply of Married Women and Spousal Tax Deductions in Japan.”} Coauthored with Minchung Hsu, Pei-Ju Liao, and Yi-Chan Tsai. 
    \begin{itemize}
        \item[] Population aging results in the problem of labor shortage. The labor force of Japanese females presents a solution to the Japanese government. Despite measures to improve the female labor force participation rate, encouraging women to work more or even shift from part-time jobs to full-time jobs was considered by the government. This study employs a standard incomplete markets model (or the Bewley-Huggett-Aiyagari style model) with heterogeneous households to examine the impact of two tax reforms (2004 and 2018) and a counterfactual reform abolishing spousal deduction on women's employment. Findings indicate that reducing spousal deduction after the 2004 tax reform boosts labor participation due to negative wealth effects, while the 2018 income threshold extensions result in a shift from full-time to part-time employment. Conversely, completely abolishing spousal deduction encourages a shift from part-time to full-time employment.
    \end{itemize}
    \item \textbf{“A Two-Stage Pseudo-Maximum Likelihood Method for Poisson Models with Two Binary Endogenous Explanatory Variables: An Application to Trade and Investment Agreements.”} Coauthored with Chingmu Chen, Shin-Kun Peng, and Wen-Jen Tsay. 
    \begin{itemize}
        \item[] This paper introduces a novel estimator for the Poisson model with dual binary endogenous explanatory variables, addressing the need to quantify the effects of preferential economic integration agreements. We develop a two-stage pseudo-maximum likelihood estimator and derive analytical gradient matrices as well as Hessian matrices. These exact formulas improve computational speed significantly. Applying our approach to trade data, we find that preferential trade agreements have a positive impact on bilateral trade flows, while bilateral investment treaties have a negative effect. Additionally, the positive interaction term between these two treaties suggests that preferential trade agreements promote horizontal foreign direct investment.
    \end{itemize}

\newpage

\section*{Working in Progress}
    \item \textbf{“How Can We Trust the Inference for the Interrupted Time Series Model?”}
    \begin{itemize}
        \item[] This paper addresses a critical gap in the statistical foundations of interrupted time series (ITS) analysis. We demonstrate that ordinary least squares (OLS) estimators in ITS models are inherently superconsistent due to deterministic time trends, converging faster than the usual square root rate and requiring proper rescaling. We propose an invertible scaling matrix and derive the asymptotic distribution of the OLS estimator for a standard single-treatment-group ITS model with serially dependent errors. Monte Carlo simulations show that our approach, accounting for superconsistency of the time trend coefficient, yields more accurate hypothesis tests than conventional methods across various autocorrelation levels.
    \end{itemize}
    \item \textbf{“Estimating Count Data Models with Endogenous Switching Variables.”}
    \begin{itemize}
        \item[] This paper proposes a novel approach to address the estimation and inference problems concerning the count data models featuring $k$ binary endogenous explanatory variables. The proposed estimator fills the gap in various economic domains where discussions involve endogenous dummy variables (treatment effects) and sample selection. Under the correct model specification, the full information maximum likelihood approach to estimation is introduced. Additionally, we develop both a one-stage quasi-maximum likelihood estimator and a two-stage quasi-maximum likelihood estimator to offer relatively robust alternatives. 
    \end{itemize}
    \item \textbf{“Asset-Based versus Income-Based Credit Constraints: Evidence from U.S. Business Cycles”}
    \begin{itemize}
        \item[] This paper assesses the impact of different credit constraints on U.S. business cycles by expanding the two-sector DSGE model of Iacoviello and Neri (2010) to include both collateral- and income-based credit constraints. Using a Bayesian approach, the study examines four scenarios: no borrowing, collateral-based constraints only, income-based constraints only, and a hybrid model that combines both. Analysis of U.S. data from 1965 to 2021 shows that the hybrid model has the highest posterior marginal log-likelihood, accurately reflecting procyclical patterns observed in the data. In addition, this model better captures the relationships among output, non-residential investment, and inflation compared to other specifications. Finally, cost-push shocks are identified as the primary source of economic fluctuations, accounting for significant variations in key economic indicators.
    \end{itemize}
\end{enumerate}



% talks
\cvsection{Conference Presentation}
\indent 

Macroeconometric Modelling Workshop -- 2023. “Financial Frictions and Uncertainty Shocks.”

Taiwan Economic Association Conference -- 2016. “Estimating the Willingness to Pay for the Civic Value When Absentee Voting Is Not Allowed.”



% teaching experience 
\cvsection{Teaching Experience}
\indent

\textbf{Macroeconomic Theory (I)}\hfill Fall 2019 -- 2022

\hspace{2em}Teaching Assistant at National Taiwan University

\textbf{Macroeconomic Theory (II)}\hfill Spring 2020 -- 2021, 2024

\hspace{2em}Teaching Assistant at National Taiwan University

\textbf{Theory of Income and Employment (III)}\hfill Fall 2022

\hspace{2em}Teaching Assistant at National Taiwan University

\textbf{Economics: Empirical Study and Forecasting (I)}\hfill Fall 2020--2022

\hspace{2em}Teaching Assistant at National Taiwan University

\textbf{Economics: Empirical Study and Forecasting (II)}\hfill Spring 2019--2021

\hspace{2em}Teaching Assistant, at National Taiwan University

\textbf{Basic Macroeconomics}\hfill Summer 2022

\hspace{2em}Teaching Assistant at National Taiwan University

\textbf{Basic Quantitative Method}\hfill Summer 2020--2022

\hspace{2em}Teaching Assistant at National Taiwan University

\textbf{Principle of Economics (I)}\hfill Fall 2019--2020

\hspace{2em}Teaching Assistant at National Taiwan University

\textbf{Principle of Economics (II)}\hfill Spring 2019--2020

\hspace{2em}Teaching Assistant at National Taiwan University


% industrial experience
% \section*{Industrial Experience}
% \indent

% \textbf{Strawberry Inc.}, Beijing, China

% \hspace{2em}Co-funder, 2077 - 

% \textbf{Delicious Mango Holdings Ltd.}, London, UK

% \hspace{2em}\LaTeX~Engineer Intern, 2076 - 2077

% \textbf{CV Agency Co. Ltd.}, Hawaii, USA

% \hspace{2em}CV Broker, 2075 - 2076
% \begin{itemize}[leftmargin=5em, topsep=0pt, itemsep=0pt]
%     \item Helped a three-year old child create her first CV
%     \item Hilight 2 ...
%     \item Hilight 3 ...
% \end{itemize}




% language
%\cvsection{Language}
%\indent

%Native Chinese (Mandarin \& Northern Shaanxi Dialect)

%Fluent English




% reference 
\cvsection{Reference}
\indent

\begin{tabular}{lrr}
% Referee 1
\begin{minipage}[t]{4.8cm}
Yi-Chan Tsai                    \\
Department of Economics         \\
National Taiwan University      \\
No. 1, Sec. 4, Roosevelt Road,  \\
Taipei 10617, Taiwan            \\
+886-2-3366-8314                \\
\href{yichantsai@ntu.edu.tw}{yichantsai\textrm{@}ntu.edu.tw}
\end{minipage}
&
% Referee 2
\begin{minipage}[t]{4.8cm}
Wen-Jen Tsay                    \\
Department of Economics         \\
Soochow University              \\
No. 56, Sec. 1, GuiYang St.,    \\
Taipei 10048, Taiwan            \\
%+886-2-2782-2791\#296           \\
\href{wtsay@econ.sinica.edu.tw}{wtsay\textrm{@}econ.sinica.edu.tw}
\end{minipage}
&
\begin{minipage}[t]{4.8cm}
Pei-Ju  Liao                    \\
Department of Economics         \\
National Taiwan University      \\
No. 1, Sec. 4, Roosevelt Road,  \\
Taipei 10617, Taiwan            \\
+886-2-3366-8371                \\
\href{mailto:pjliao@ntu.edu.tw}{pjliao\textrm{@}ntu.edu.tw}
\end{minipage}
\end{tabular}


%\vfill

%\section*{\hfill\color{OliveGreen}\today}

\end{document}
